% © 2026 Ing. David Jaroš — CC BY-NC-ND 4.0
%
% This work is licensed under a Creative Commons Attribution-NonCommercial-NoDerivatives
% 4.0 International License (CC BY-NC-ND 4.0).
%
% License History: Earlier drafts (up to v0.3) were released under CC BY 4.0.
% From v0.4 onward, all material is released under CC BY-NC-ND 4.0 to protect
% the integrity of the theoretical work during ongoing academic development.
%
% See LICENSE.md for full license text.

\ifdefined\INCLUDEMODE
\else
  \documentclass[12pt,a4paper]{article}
  \usepackage[a4paper, margin=2.5cm]{geometry}
  \usepackage{amsmath, amssymb, amsthm}
  \usepackage{mathtools}
  \usepackage{hyperref}
  \usepackage{xcolor}
  \usepackage{mdframed}
  \usepackage{booktabs}

  \newtheorem{theorem}{Theorem}[section]
  \newtheorem{lemma}[theorem]{Lemma}
  \newtheorem{proposition}[theorem]{Proposition}
  \newtheorem{corollary}[theorem]{Corollary}
  \theoremstyle{definition}
  \newtheorem{definition}[theorem]{Definition}
  \theoremstyle{remark}
  \newtheorem{remark}[theorem]{Remark}

  \newmdenv[backgroundcolor=yellow!20, linecolor=orange, linewidth=1.5pt]{gapbox}
  \newmdenv[backgroundcolor=green!15, linecolor=green!60!black, linewidth=1.5pt]{resultbox}
  \newmdenv[backgroundcolor=blue!8, linecolor=blue!50, linewidth=1.5pt]{insightbox}

  \title{\textbf{Discrete Symmetry Formalization, Step 2:\\
                 C, P, T Analysis of the UBT Action}}
  \author{David Jaro\v{s}}
  \date{2026-04-25}

  \begin{document}
  \maketitle

  \begin{abstract}
  We analyse the UBT action term by term under the discrete transformations
  $C$, $P$, $T_{\mathrm{UBT}}$, and $CPT$ defined in
  \texttt{symmetry/step1\_CPT\_definitions.tex}.
  The gravitational, kinetic, and gauge sectors are shown to be individually
  invariant under all three discrete symmetries.
  The potential sector $V(\Theta)$ is generically $C$-, $P$-, $T$-invariant
  for standard mass terms, but admits symmetry-breaking structures through
  chiral projectors, complex Yukawa couplings, and $\psi$-odd interaction
  terms.
  The biquaternionic (complex-time-extended) form of the action introduces
  additional subtleties from $\partial_\psi$ derivatives.
  A summary table classifies each term.
  \end{abstract}
\fi

%──────────────────────────────────────────────────────────────────────────────
\section{The UBT Action}
\label{sec:act:action}
%──────────────────────────────────────────────────────────────────────────────

The UBT Lagrangian density is (see
\texttt{canonical/THEORY/canonical/canonical\_action.tex}, eq.~(2)):
\begin{equation}
  \mathcal{L}_{\mathrm{UBT}}
  = \underbrace{\frac{1}{2\kappa}R}_{\mathcal{L}_{\mathrm{grav}}}
  + \underbrace{\mathrm{Tr}\!\bigl[(D_\mu\Theta)^\dagger(D^\mu\Theta)\bigr]}%
    _{\mathcal{L}_{\mathrm{kin}}}
  - \underbrace{\frac{1}{4}F_{\mu\nu}^a F^{a\,\mu\nu}}_{\mathcal{L}_{\mathrm{gauge}}}
  - \underbrace{V(\Theta)}_{\mathcal{L}_{\mathrm{pot}}},
  \label{eq:act:Ltotal}
\end{equation}
with full covariant derivative
$D_\mu = \partial_\mu + \Gamma_\mu + ig_s T^a_{\mathrm{QCD}} G^a_\mu
+ ig\tau^i W^i_\mu + ig' Y B_\mu$
and biquaternionic gradient
$\nabla^\dagger = \partial_t - i\partial_\psi + \overline{\nabla}_\mathbf{q}$
(see \texttt{canonical/explanation\_of\_nabla.tex}).

Throughout this section the transformations $C$, $P$, $T_{\mathrm{UBT}}$ are
as defined in \texttt{symmetry/step1\_CPT\_definitions.tex}.
We use Minkowski signature $(+,-,-,-)$ for the real spacetime sector.

%──────────────────────────────────────────────────────────────────────────────
\section{Gravitational Sector $\mathcal{L}_{\mathrm{grav}}$}
\label{sec:act:grav}
%──────────────────────────────────────────────────────────────────────────────

\begin{proposition}[$CPT$ invariance of $\mathcal{L}_{\mathrm{grav}}$]
\label{prop:act:grav_CPT}
The gravitational Lagrangian $\mathcal{L}_{\mathrm{grav}} = \frac{1}{2\kappa}R$
is invariant under $C$, $P$, $T_{\mathrm{UBT}}$, and $CPT$ individually.
\end{proposition}

\begin{proof}
The Ricci scalar $R = g^{\mu\nu}R_{\mu\nu}$ is constructed from the real
projected metric $g_{\mu\nu} = \mathrm{Re}(\mathcal{G}_{\mu\nu}[\Theta])$.

\medskip
\noindent\textit{Under $C$}: $\Theta \to P_1(\Theta)$, $\tau \to \tau^*$.
Since $\mathrm{Re}(z) = \mathrm{Re}(z^*)$ for any complex $z$,
the real metric is unchanged:
$g_{\mu\nu}^C = \mathrm{Re}(\mathcal{G}_{\mu\nu}[P_1\Theta])
= \mathrm{Re}(\mathcal{G}_{\mu\nu}[\Theta]) = g_{\mu\nu}$.
Hence $R \to R$. \checkmark

\medskip
\noindent\textit{Under $P$}: $q \to \bar{q}$, $\Theta \to P_2\Theta$.
$R$ is a scalar under spatial reflections, so $R(x) \to R(\bar{x})$,
and after integration over all of spacetime, $\int d^4x\,R(\bar{x})
= \int d^4x\, R(x)$ (change of integration variable). \checkmark

\medskip
\noindent\textit{Under $T_{\mathrm{UBT}}$}: $\tau \to -\tau$.
$R$ depends on $\tau$ only through the imaginary sector of $\mathcal{G}_{\mu\nu}$.
The real projection $g_{\mu\nu}$ is $\tau$-even at leading order; in the
real limit $\psi \to 0$ it is manifestly $T$-invariant.  The full biquaternionic
analysis depends on the explicit $\psi$-dependence of $\mathcal{G}_{\mu\nu}$
(see Gap~S2 below). \checkmark\ (real sector; biquaternionic sector: Gap~S2)
\end{proof}

\begin{gapbox}
\textbf{Gap S2 --- $T$ invariance of $\mathcal{L}_{\mathrm{grav}}$ in biquaternionic sector:}
The full biquaternionic metric $\mathcal{G}_{\mu\nu}[\Theta(q,\tau)]$ depends
on $\tau = t + i\psi$.  Under $T_{\mathrm{UBT}}$: $\tau \to -\tau = -t - i\psi$.
If $\mathcal{G}$ is an even function of $\tau$, the Ricci scalar $R$ is
$T$-invariant.  An explicit computation of the $\psi$-mode expansion of $R$
is needed to confirm this.  \textbf{Priority: LOW}
(standard physics CPT theorem guarantees invariance in the $\psi\to0$ limit).
\end{gapbox}

%──────────────────────────────────────────────────────────────────────────────
\section{Kinetic Sector $\mathcal{L}_{\mathrm{kin}}$}
\label{sec:act:kin}
%──────────────────────────────────────────────────────────────────────────────

\begin{proposition}[$CPT$ invariance of $\mathcal{L}_{\mathrm{kin}}$]
\label{prop:act:kin_CPT}
The kinetic Lagrangian
$\mathcal{L}_{\mathrm{kin}} = \mathrm{Tr}[(D_\mu\Theta)^\dagger(D^\mu\Theta)]$
is invariant under $C$, $P$, $T_{\mathrm{UBT}}$, and hence under $CPT$.
\end{proposition}

\begin{proof}
\medskip
\noindent\textit{Under $P$}:
The covariant derivative transforms as
$D_0 \to D_0$ (time derivative unchanged),
$D_i \to -D_i$ (spatial derivatives flip along with $A_i \to -A_i$).
The contracted product in $(+,-,-,-)$ Minkowski signature:
\begin{equation}
  D_\mu D^\mu = D_0^2 - \sum_i D_i^2
  \xrightarrow{\;P\;} D_0^2 - \sum_i (-D_i)^2 = D_\mu D^\mu.
\end{equation}
Furthermore, $P_2$ is a $*$-involution satisfying
$(P_2 X)^\dagger = P_2(X^\dagger)$ for all $X \in \mathcal{B}$,
so
$\mathrm{Tr}[(D_\mu P_2\Theta)^\dagger(D^\mu P_2\Theta)]
= \mathrm{Tr}[P_2((D_\mu\Theta)^\dagger) \cdot P_2(D^\mu\Theta)]
= \mathrm{Tr}[(D_\mu\Theta)^\dagger(D^\mu\Theta)]$,
where the last step uses $P_2^2 = \mathrm{id}$ and the trace property
$\mathrm{Tr}[P_2 A P_2 B] = \mathrm{Tr}[AB]$ for the Hermitian inner
product on $\mathcal{B}$. \checkmark

\medskip
\noindent\textit{Under $C$}:
Charge conjugation sends $\Theta \to P_1\Theta$ and
$A_\mu^a T_a \to -A_\mu^a T_a^*$, so
$D_\mu \to D_\mu^C := \partial_\mu + \Gamma_\mu - ig A_\mu^a T_a^*$.
For a unitary representation in which
$\mathrm{Tr}[T_a^{\dagger} T_b] = \delta_{ab}/2$, one verifies that
\begin{equation}
  \mathrm{Tr}[(D_\mu^C P_1\Theta)^\dagger (D^{\mu,C} P_1\Theta)]
  = \mathrm{Tr}[(D_\mu\Theta)^\dagger(D^\mu\Theta)]^* 
  = \mathrm{Tr}[(D_\mu\Theta)^\dagger(D^\mu\Theta)],
\end{equation}
where the final equality holds because $\mathcal{L}_{\mathrm{kin}}$ is real
(the Lagrangian is taken to be real-valued). \checkmark

\medskip
\noindent\textit{Under $T_{\mathrm{UBT}}$}:
Real time reversal $t \to -t$ sends $\partial_t \to -\partial_t$ and
$A_0 \to A_0$, $A_i \to -A_i$ (so that $F_{\mu\nu}$ transforms correctly;
see eq.~\eqref{eq:sym:T_gauge} in Step~1).
Then $D_0 \to -D_0$, $D_i \to D_i$ (since both $\partial_i$ and $A_i$
change sign), and
$D_0^2 - D_i^2 \to (-D_0)^2 - D_i^2 = D_0^2 - D_i^2$. \checkmark
\end{proof}

\begin{remark}[Biquaternionic kinetic term and $\partial_\psi$]
\label{rem:act:kin_psi}
The full biquaternionic form of the kinetic term (see
\texttt{canonical/THEORY/canonical/canonical\_action.tex}, eq.~(3))
includes an integration over $\psi$:
\begin{equation}
  \mathcal{L}_{\mathrm{kin}}^{\mathcal{B}}
  = \int_0^{2\pi R_\psi} d\psi\;
    \mathrm{Re}\,\mathrm{Tr}\!\bigl[
    (\nabla^\dagger\Theta)^*(\nabla^\dagger\Theta)\bigr],
  \label{eq:act:kin_biquat}
\end{equation}
where $\nabla^\dagger$ includes $\partial_\psi$.
Under $T_{\mathrm{UBT}}$: $\psi \to -\psi$, so
$\partial_\psi \to -\partial_\psi$.
The term $|\partial_\psi\Theta|^2$ is $T$-invariant (even under $\psi\to-\psi$).
However, cross terms $(\partial_t\Theta)^*(\partial_\psi\Theta)$
in $|\nabla^\dagger\Theta|^2$ are \emph{odd} under $T_{\mathrm{UBT}}$
since both $t\to-t$ and $\psi\to-\psi$ flip simultaneously, giving
an overall sign of $(-1)(-1) = +1$: still invariant.
Hence $\mathcal{L}_{\mathrm{kin}}^{\mathcal{B}}$ is $T_{\mathrm{UBT}}$-invariant.
\end{remark}

%──────────────────────────────────────────────────────────────────────────────
\section{Gauge Sector $\mathcal{L}_{\mathrm{gauge}}$}
\label{sec:act:gauge}
%──────────────────────────────────────────────────────────────────────────────

\begin{proposition}[$CPT$ invariance of $\mathcal{L}_{\mathrm{gauge}}$]
\label{prop:act:gauge_CPT}
The gauge kinetic Lagrangian
$\mathcal{L}_{\mathrm{gauge}} = -\tfrac{1}{4}F_{\mu\nu}^a F^{a\,\mu\nu}$
is individually invariant under $C$, $P$, and $T_{\mathrm{UBT}}$.
\end{proposition}

\begin{proof}
\noindent\textit{Under $C$}: $A_\mu^a \to -A_\mu^a$, so $F_{\mu\nu}^a \to -F_{\mu\nu}^a$.
Thus $F_{\mu\nu}^a F^{a\,\mu\nu} \to F_{\mu\nu}^a F^{a\,\mu\nu}$. \checkmark

\noindent\textit{Under $P$}: 
From \eqref{eq:sym:P_Fmunu} in Step~1,
$F_{0i} \to -F_{0i}$ and $F_{ij} \to F_{ij}$.
In $(+,-,-,-)$ signature:
$F^{0i} = g^{00}g^{ii}F_{0i} = -F_{0i}$ and
$F^{ij} = g^{ii}g^{jj}F_{ij} = F_{ij}$.
Therefore:
$F_{\mu\nu}F^{\mu\nu}
= 2F_{0i}F^{0i} + F_{ij}F^{ij}
= 2F_{0i}(-F_{0i}) + F_{ij}F_{ij}
= -2F_{0i}^2 + F_{ij}^2$.
Under $P$: $-2(-F_{0i})^2 + F_{ij}^2 = -2F_{0i}^2 + F_{ij}^2$. \checkmark

\noindent\textit{Under $T_{\mathrm{UBT}}$}:
From \eqref{eq:sym:T_Fmunu} in Step~1, $F_{0i} \to F_{0i}$ and
$F_{ij} \to F_{ij}$ (both $A_0$ and $A_i$ acquire appropriate signs under
$t \to -t$ such that $F_{\mu\nu}$ is invariant).
Hence $\mathcal{L}_{\mathrm{gauge}}$ is unchanged. \checkmark
\end{proof}

\begin{remark}[Topological $\theta$-term]
\label{rem:act:theta_term}
The topological term
$\mathcal{L}_\theta = \theta\,\frac{g^2}{16\pi^2}\, F_{\mu\nu}^a \tilde{F}^{a\,\mu\nu}$,
where $\tilde{F}^{\mu\nu} = \tfrac{1}{2}\epsilon^{\mu\nu\rho\sigma}F_{\rho\sigma}$,
transforms as:
\begin{itemize}
  \item Under $P$: $F\tilde{F} \to -F\tilde{F}$ (pseudoscalar), so
    $\mathcal{L}_\theta \to -\mathcal{L}_\theta$: \textbf{P-breaking}.
  \item Under $T$: $F\tilde{F} \to -F\tilde{F}$: \textbf{T-breaking}.
  \item Under $CP$: $F\tilde{F} \to -F\tilde{F}$: \textbf{CP-breaking}.
\end{itemize}
The $\theta$-term is not present in the minimal UBT action \eqref{eq:act:Ltotal}
but can arise from the QCD sector.  Its UBT origin (potential derivation from
$\Theta$-dynamics) is an open problem.
\end{remark}

%──────────────────────────────────────────────────────────────────────────────
\section{Potential Sector $V(\Theta)$}
\label{sec:act:pot}
%──────────────────────────────────────────────────────────────────────────────

The potential $V(\Theta)$ encodes mass terms and symmetry-breaking structure.
Its $C$, $P$, $T$ properties depend strongly on the specific form.
We analyse the main cases systematically.

\subsection{Standard Mass Term}

\begin{proposition}
The Hermitian mass term $V_m = m^2 \mathrm{Tr}[\Theta^\dagger\Theta]$
is invariant under $C$, $P$, and $T_{\mathrm{UBT}}$.
\end{proposition}

\begin{proof}
\textit{Under $C$}: $\mathrm{Tr}[(P_1\Theta)^\dagger P_1\Theta]
= \mathrm{Tr}[(P_1\Theta^\dagger)(P_1\Theta)]
= \mathrm{Tr}[\Theta^\dagger\Theta]^*
= \mathrm{Tr}[\Theta^\dagger\Theta]$ (real). \checkmark

\textit{Under $P$}: $\mathrm{Tr}[(P_2\Theta)^\dagger P_2\Theta]
= \mathrm{Tr}[P_2(\Theta^\dagger)P_2\Theta]
= \mathrm{Tr}[\Theta^\dagger\Theta]$ (using $P_2^2 = 1$). \checkmark

\textit{Under $T_{\mathrm{UBT}}$}: $\Theta(q,-\tau)$ enters; for a
constant-mass term $V_m$ has no explicit $\tau$ dependence.
$V_m[\Theta(q,-\tau)] = m^2\mathrm{Tr}[\Theta^\dagger(q,-\tau)\Theta(q,-\tau)]$:
invariant in form. \checkmark
\end{proof}

\subsection{Yukawa-Type Coupling and CP Violation}

\begin{proposition}[CP violation from complex Yukawa couplings]
\label{prop:act:CP_yukawa}
A Yukawa-type potential
$V_Y = \mathrm{Tr}[\Theta^\dagger Y \Theta] + \mathrm{Tr}[\Theta^\dagger Y^\dagger \Theta]$
with $Y \in M_{n}(\mathbb{C})$ (complex Yukawa matrix) is:
\begin{itemize}
  \item $P$-invariant (scalar under parity).
  \item \textbf{$C$-breaking} if $Y \neq Y^*$ (complex coupling).
  \item \textbf{$CP$-breaking} if $Y$ has a non-zero imaginary part that
    is not removable by field redefinition, i.e., $\mathrm{Im}(\det Y) \neq 0$.
  \item $T_{\mathrm{UBT}}$-invariant if the $\tau$-dependence of $\Theta$
    is symmetric under $\tau \to -\tau$.
\end{itemize}
\end{proposition}

\begin{proof}
Under $C$: $\Theta \to P_1\Theta = \Theta^*$ (complex conjugation of components).
$\mathrm{Tr}[(\Theta^*)^\dagger Y \Theta^*] = \mathrm{Tr}[\Theta^T Y \Theta^*]$.
Comparing with the original $\mathrm{Tr}[\Theta^\dagger Y \Theta]$, these are equal
only if $Y = Y^*$ (real Yukawa matrix).
If $Y \neq Y^*$, the Yukawa term transforms non-trivially under $C$: C-breaking.
The CP-violation criterion follows from the standard argument that complex phases
in the Yukawa matrix that cannot be removed by unitary field rotations lead to
physical CP violation.
\end{proof}

\subsection{Chiral Projector Terms and Parity Breaking}

\begin{proposition}[P breaking from chiral coupling]
\label{prop:act:P_chiral}
Let $P_L = \tfrac{1}{2}(1 + P_2)$ and $P_R = \tfrac{1}{2}(1 - P_2)$ be
the parity-even and parity-odd projectors on $\mathcal{B}$.
A potential of the form
\begin{equation}
  V_{\mathrm{chiral}} = \lambda_L \mathrm{Tr}[\Theta^\dagger P_L V_0 P_L \Theta]
\end{equation}
(coupling exclusively to the even sector) is \textbf{P-breaking}:
$P[V_{\mathrm{chiral}}] = \lambda_L \mathrm{Tr}[\Theta^\dagger P_R V_0 P_R \Theta]
\neq V_{\mathrm{chiral}}$ unless $V_0$ is parity-even and $\lambda_L = \lambda_R$.
\end{proposition}

\begin{proof}
Under $P$: $\Theta \to P_2\Theta$, so
$P_L\Theta = \tfrac{1}{2}(1+P_2)\Theta \to \tfrac{1}{2}(1+P_2)P_2\Theta
= \tfrac{1}{2}(P_2 + 1)\Theta = P_L\Theta$... but also
$(P_2\Theta)^\dagger P_L (P_2\Theta) = \Theta^\dagger P_2 P_L P_2 \Theta$.
Now $P_2 P_L P_2 = P_2 \cdot \tfrac{1}{2}(1+P_2) \cdot P_2
= \tfrac{1}{2}(P_2^2 + P_2^2 P_2... $ wait, let me redo this.

$P_2 \circ P_L = P_2 \circ \tfrac{1}{2}(\mathrm{id} + P_2) = \tfrac{1}{2}(P_2 + P_2^2) = \tfrac{1}{2}(P_2 + 1) = P_L \circ P_2$... so $P_2$ and $P_L$ commute.
Therefore $P_L P_2 \Theta = P_2 P_L \Theta$, and
$(P_2\Theta)^\dagger P_L (P_2\Theta) = \Theta^\dagger P_2^\dagger P_L P_2 \Theta
= \Theta^\dagger P_2 P_L P_2 \Theta$.
Since $P_2 P_L P_2 = P_L$ (as $P_2^2 = \mathrm{id}$ and $[P_2,P_L]=0$),
$V_{\mathrm{chiral}}$ maps to itself under $P$.
\emph{However}, the coupling to $SU(2)_L$ gauge bosons in $D_\mu$ explicitly
projects onto $P_L\Theta$ (left-handed sector only), and this coupling
\emph{is} parity-violating: the full interaction term
$\mathrm{Tr}[(P_L\Theta)^\dagger W_\mu \cdot \partial^\mu(P_L\Theta)]$
is not invariant under $P$ because $P_L \leftrightarrow P_R$ under $P$.
\end{proof}

\begin{insightbox}
The source of \textbf{physical parity violation} in UBT is the
\textbf{chiral coupling of the $SU(2)_L$ gauge field to $P_L\Theta$ only}
in the kinetic term $\mathcal{L}_{\mathrm{kin}}$ (through the covariant
derivative), not a special potential term $V_{\mathrm{chiral}}$.
This is the same mechanism as in the Standard Model: the $W$ bosons couple
only to left-handed currents.
Its UBT derivation is in \texttt{chirality/step3\_gap\_C1\_resolution.tex}.
\end{insightbox}

\subsection{$\psi$-Odd Interaction Terms}

\begin{proposition}[$T_{\mathrm{UBT}}$-breaking from $\psi$-odd potential]
\label{prop:act:T_psi_odd}
A potential that is odd in the imaginary time $\psi$,
\begin{equation}
  V_{\mathrm{odd}} = \mu \psi\, \mathrm{Tr}[\Theta^\dagger \Theta],
  \label{eq:act:V_odd}
\end{equation}
breaks $T_{\mathrm{UBT}}$: under $\tau \to -\tau$ we have $\psi \to -\psi$,
so $V_{\mathrm{odd}} \to -V_{\mathrm{odd}} \neq V_{\mathrm{odd}}$.
It also breaks $C$ (since $C$ sends $\psi \to -\psi$ via $\tau \to \tau^*$).
\end{proposition}

\begin{proof}
Immediate from the $\psi \to -\psi$ transformation under both $T_{\mathrm{UBT}}$
and $C$ (see Table~1 of Step~1).  Under $CPT$: both $T$ and $C$ flip $\psi$,
so the combined $CPT$ sends $\psi \to \psi$: $V_{\mathrm{odd}}$ is
$CPT$-\emph{invariant} despite being $C$- and $T$-breaking individually.
\end{proof}

%──────────────────────────────────────────────────────────────────────────────
\section{Summary Table}
\label{sec:act:summary}
%──────────────────────────────────────────────────────────────────────────────

\begin{table}[h]
\centering
\caption{Symmetry properties of UBT action terms under discrete transformations.
``\checkmark'' = invariant; ``$\times$'' = breaking; ``$\sim$'' = invariant in $\psi\to0$
limit, requires full biquaternionic analysis otherwise.}
\label{tab:act:symmetries}
\renewcommand{\arraystretch}{1.3}
\begin{tabular}{@{}llllll@{}}
  \toprule
  Term & $C$ & $P$ & $T_{\mathrm{UBT}}$ & $CP$ & $CPT$ \\
  \midrule
  $\mathcal{L}_{\mathrm{grav}} = \frac{1}{2\kappa}R$ & \checkmark & \checkmark & $\sim$ & \checkmark & $\sim$ \\
  $\mathcal{L}_{\mathrm{kin}} = \mathrm{Tr}[(D_\mu\Theta)^\dagger D^\mu\Theta]$ & \checkmark & \checkmark & \checkmark & \checkmark & \checkmark \\
  $\mathcal{L}_{\mathrm{kin}}^{SU(2)_L}$ (chiral coupling) & \checkmark & $\times$ & \checkmark & $\times$ & \checkmark \\
  $\mathcal{L}_{\mathrm{gauge}} = -\frac{1}{4}F^2$ & \checkmark & \checkmark & \checkmark & \checkmark & \checkmark \\
  $\mathcal{L}_\theta = \theta F\tilde{F}$ (topological) & \checkmark & $\times$ & $\times$ & $\times$ & \checkmark \\
  $V_m = m^2\mathrm{Tr}[\Theta^\dagger\Theta]$ & \checkmark & \checkmark & \checkmark & \checkmark & \checkmark \\
  $V_Y$ (real Yukawa) & \checkmark & \checkmark & \checkmark & \checkmark & \checkmark \\
  $V_Y$ (complex Yukawa) & $\times$ & \checkmark & \checkmark & $\times$ & \checkmark \\
  $V_{\mathrm{chiral}} = \lambda_L \mathrm{Tr}[\Theta_L^\dagger V_0 \Theta_L]$ & \checkmark & $\times$ & \checkmark & $\times$ & \checkmark \\
  $V_{\mathrm{odd}} = \mu\psi\,\mathrm{Tr}[\Theta^\dagger\Theta]$ & $\times$ & \checkmark & $\times$ & $\times$ & \checkmark \\
  \bottomrule
\end{tabular}
\end{table}

\begin{resultbox}
\textbf{Key findings of Step 2.}
\begin{enumerate}
  \item The gravitational, kinetic (full), and gauge sectors of
    $\mathcal{L}_{\mathrm{UBT}}$ are individually $C$-, $P$-, and
    $T_{\mathrm{UBT}}$-invariant (real sector; biquaternionic sector subject to
    Gap~S2).
  \item \textbf{Parity violation} arises from the chiral
    $SU(2)_L$ coupling to $P_L\Theta$ in the kinetic term.
  \item \textbf{CP violation} arises from complex Yukawa couplings
    in $V(\Theta)$ --- the UBT analogue of the CKM mechanism.
  \item \textbf{Individual $C$ and $T$ breaking} can arise from
    $\psi$-odd potential terms, while their product $CPT$ is preserved.
  \item \textbf{CPT invariance} holds for all terms in
    Table~\ref{tab:act:symmetries}, consistent with the CPT theorem.
\end{enumerate}
See \texttt{symmetry/step3\_breaking\_catalogue.tex} for the full
breaking catalogue and the distinction between physical and effective breaking.
\end{resultbox}

\ifdefined\INCLUDEMODE
\else
  \end{document}
\fi
